% GENERATED FILE. Edit canonical lessons, then run npm run generate:books.
% canonical-source-hash: fnv1a64:13ab5b3394e0abc0
% canonical-lessons: ES-C474-invitar, ES-C474-importe, ES-C474-descontar, ES-C474-golpe, ES-R474-repaso-pago, ES-C474-sintesis-pago

\chapter{El curso que paga la empresa}
\label{ch:es-pago}
\hlchaptermodality{474}

\begin{chapteropening}
\textbf{By the end of this chapter:} \emph{I can follow a conversation about who pays for something: hear that an amount is deducted over months rather than taken all at once, and match that to a statement saying the opposite in different words.}

Match a statement about paying the full amount at once against a speaker who says the opposite in different words.
\end{chapteropening}

\section[invitar]{invitar — to invite}

\label{lesson:ES-C474-invitar}



\begin{quote}
Say \textbf{el invitado} and \textbf{pagar}. Today's verb made the first, and in one very common use it means the second.
\end{quote}

\subsection*{You'll want to know: invitar}
\textbf{invitar} — to invite.

\begin{quote}
Nos \textbf{invitaron} a la fiesta.
\end{quote}

You have met this verb once without being handed it. El invitado named it in passing — \emph{\textquotedblleft{}from \textbf{invitar}, Latin \textbf{invitare}, to invite; \textbf{invitado} is its participle\textquotedblright{}} — and then stayed with the guest. Here it is as yours.

\emph{Invitar a alguien a algo} — the person with \textbf{a}, and the thing with \textbf{a} as well.

\begin{grammarlens}[title={Grammar Lens: in Spanish, inviting can mean paying}]
This is the use you will meet first in real life, and English has no single word for it.

\noindent
\begin{tabularx}{\linewidth}{@{}>{\raggedright\arraybackslash}X>{\raggedright\arraybackslash}X@{}}
\toprule
\textbf{Spanish} & \textbf{what it means} \\
\midrule
\emph{te invito} & \textbf{I'm paying} \\
\emph{invito yo} & this one's on me \\
\bottomrule
\end{tabularx}

Somebody who says \emph{te invito a un café} is not asking whether you are free. They are telling you they will pay for it. Answering \emph{vale, ¿a qué hora?} when the coffee is in front of you will land oddly.

So \emph{invitar} covers both the asking and the paying, and which one is meant is settled by whether a time is mentioned.
\end{grammarlens}

\begin{cousinweb}
Latin \textbf{invitare}, to invite — and the honest answer about what is inside it is that nobody is sure. The older guess tied it to a root about \textbf{willingness}, which would make an invitation an appeal to somebody's good will.

English took the whole word twice: \textbf{invite}, and the noun \textbf{invitation}. There is no native English word for this, which is a small clue that the custom travelled with the Latin.

What Spanish added is the paying. That is not in the Latin and not in the English — it is a courtesy that attached itself to the word here.
\end{cousinweb}

\subsection*{Your turn}
Say these aloud:

\begin{itemize}
\raggedright
  \item they invited us to the party
  \item this one's on me
  \item the noun this verb already gave you
\end{itemize}

\begin{tcolorbox}[breakable,colback=teal!4,colframe=teal!35!black,title={Before you move on}]
To invite? (\textbf{Invitar}.) What does \emph{te invito} usually mean? (\textbf{I'm paying}.) Which word did you have first? (\textbf{El invitado} — its participle.)
\end{tcolorbox}
\section[el importe]{el importe — the amount}

\label{lesson:ES-C474-importe}



\begin{quote}
Say \textbf{importar} and \textbf{importante}. Today's noun is the third member of that family, and it is the one that appears on a bill.
\end{quote}

\subsection*{You'll want to know: el importe}
\textbf{el importe} — the amount, the sum.

\begin{quote}
Hay que pagar el \textbf{importe} completo.
\end{quote}

\emph{El importe total} — the total. \emph{El importe del curso} — the fee for a class. This is a form-and-invoice word: in speech people say \emph{cuánto es}, but anything printed says \emph{importe}.

\begin{grammarlens}[title={Grammar Lens: one verb, three jobs}]
Importar told you two of these. Here is the third, and seeing them together is what makes the word stick.

\noindent
\begin{tabularx}{\linewidth}{@{}>{\raggedright\arraybackslash}X>{\raggedright\arraybackslash}X@{}}
\toprule
\textbf{the sense} & \textbf{example} \\
\midrule
to matter & \emph{no me importa} \\
to import goods & \emph{importar café} \\
\textbf{to amount to} & \textbf{el importe} \\
\bottomrule
\end{tabularx}

The thread is Latin \textbf{in-} plus \textbf{portare}, to \textbf{carry in}. Goods carried in are imports. A thing that carries weight \emph{matters}. And what a bill \textbf{carries in} for you to pay is its \emph{importe}.

So the noun is not a fourth meaning bolted on. It is the same carrying, applied to money.
\end{grammarlens}

\begin{cousinweb}
\emph{Portare} is the carrying verb, and once you have it a great many words open at once. Spanish has already given you two of its children without saying so:

\begin{itemize}
\raggedright
  \item \textbf{el deporte} — \emph{dis-} plus \emph{portare}, carrying yourself away from work;
  \item \textbf{el comportamiento} — carrying yourself \textbf{with}, which is behaviour.
\end{itemize}

English carries the same root in \textbf{port}, \textbf{portable}, \textbf{export}, \textbf{transport} and \textbf{report} — a thing carried back.

And \emph{importante} is literally \textbf{carrying-in-ish}: something that brings weight with it.
\end{cousinweb}

\subsection*{Your turn}
Say these aloud:

\begin{itemize}
\raggedright
  \item you have to pay the full amount
  \item the three jobs of \emph{importar}
  \item what \emph{portare} means
\end{itemize}

\begin{tcolorbox}[breakable,colback=teal!4,colframe=teal!35!black,title={Before you move on}]
The amount? (\textbf{El importe}.) Which verb is it from? (\textbf{Importar} — \emph{in} plus \emph{portare}.) Would you say it out loud in a shop? (\textbf{Usually not — it is a printed word}.)
\end{tcolorbox}
\section[descontar]{descontar — to deduct}

\label{lesson:ES-C474-descontar}



\begin{quote}
Say \textbf{contar} and \textbf{el importe}. Put \emph{des-} on the first and you have what happens to the second.
\end{quote}

\subsection*{You'll want to know: descontar}
\textbf{descontar} — to deduct, to take off.

\begin{quote}
Te lo \textbf{descuentan} en doce meses.
\end{quote}

\emph{El descuento} is the noun — a discount. \emph{Con descuento} — at a reduced price.

Note the stem: \emph{desc\textbf{ue}nto}, \emph{desc\textbf{ue}ntan}. The \textbf{o} breaks into \textbf{ue} under stress, exactly as it does in \emph{contar} itself — \emph{cuento}, \emph{cuentas}.

\begin{grammarlens}[title={Grammar Lens: the des- that reverses, and the one that does not}]
Cubrir set this out plainly: \emph{des-} sometimes undoes the verb and sometimes does nothing of the kind, and only the whole word tells you which.

\noindent
\begin{tabularx}{\linewidth}{@{}>{\raggedright\arraybackslash}X>{\raggedright\arraybackslash}X@{}}
\toprule
\textbf{word} & \textbf{what des- does} \\
\midrule
\textbf{descubrir} & takes the cover off — it reverses \\
\emph{disfrutar} & nothing — the prefix does not reverse \\
\textbf{descontar} & \textbf{takes off the count — it reverses} \\
\bottomrule
\end{tabularx}

This one is in the first camp, and cleanly. \emph{Contar} adds up; \emph{descontar} takes away from the total. If you can count the thing, \emph{des-} will subtract it.

\emph{Descansar} is the same family — \emph{des-} undoing \emph{cansar}, to tire — and it is where you first met the prefix doing real work.
\end{grammarlens}

\begin{cousinweb}
\emph{Contar} is Latin \textbf{computare}: \textbf{com-}, together, plus \textbf{putare}, to reckon. So counting is a \textbf{reckoning-together}, and \emph{descontar} is a reckoning taken back out.

That \emph{computare} is where English \textbf{compute} and \textbf{count} both come from — the second worn down through French — so \emph{contar} and \emph{count} are the same word wearing different amounts of Latin.

And \emph{putare} first meant to \textbf{prune} a tree. Reckoning, in Latin, began as cutting away what does not belong, which makes \emph{descontar} an oddly exact return to the root.
\end{cousinweb}

\subsection*{Your turn}
Say these aloud:

\begin{itemize}
\raggedright
  \item they deduct it from you over twelve months
  \item the noun, meaning a discount
  \item which camp this \emph{des-} is in, and why
\end{itemize}

\begin{tcolorbox}[breakable,colback=teal!4,colframe=teal!35!black,title={Before you move on}]
To deduct? (\textbf{Descontar}.) They deduct? (\textbf{Descuentan} — the o breaks.) Does this \emph{des-} reverse? (\textbf{Yes — it takes off the count}.)
\end{tcolorbox}
\section[el golpe]{el golpe — the blow}

\label{lesson:ES-C474-golpe}



\begin{quote}
Say \textbf{descontar} and \textbf{pagar}. If the money is not taken off in instalments, it comes in one of today's word.
\end{quote}

\subsection*{You'll want to know: el golpe}
\textbf{el golpe} — the blow, the knock.

\begin{quote}
No te lo cobran de \textbf{golpe}.
\end{quote}

The literal sense is a hit: \emph{un golpe en la puerta}, a knock at the door.

But the phrase is what you need: \textbf{de golpe} — all at once, suddenly.

\begin{grammarlens}[title={Grammar Lens: the phrase a payment plan is defined against}]
Spanish has two ways to say \emph{all at once}, and a notice about money will use one to rule out the other.

\noindent
\begin{tabularx}{\linewidth}{@{}>{\raggedright\arraybackslash}X>{\raggedright\arraybackslash}X@{}}
\toprule
\textbf{phrase} & \textbf{what it pictures} \\
\midrule
\textbf{de golpe} & in one blow — suddenly, often unwelcome \\
\emph{de una vez} & in a single go — neutral \\
\bottomrule
\end{tabularx}

\emph{De golpe} carries the impact. \emph{Se me olvidó de golpe} — it went out of my head all at once. \emph{Subió de golpe} — it shot up.

So \emph{no de golpe} about a fee is doing more than saying \emph{in instalments}: it is saying you will not be hit with it. That is why the reassurance lands, and why a question about it may use the calmer \emph{de una vez} instead.
\end{grammarlens}

\begin{cousinweb}
Latin \textbf{colaphus}, a blow with the fist, borrowed from Greek \textbf{kólaphos}, a box on the ear.

The road from \emph{colaphus} to \emph{golpe} is worth a look, because two things happened: the \textbf{c} softened to \textbf{g}, and the middle syllable fell out entirely — \emph{co-la-phus} worn down to \emph{gol-pe}. That is ordinary wear, and it left a word that no longer looks Greek at all.

English has the same blow in \textbf{coup}, taken from French rather than Spanish and kept intact: a \emph{coup} is one sharp strike, which is what a \emph{golpe} is. Spanish says \textbf{golpe de estado} where English says \emph{coup d'état} — the same two words, twice.

French went further and made a quantity out of it. \textbf{Beaucoup} is \emph{beau} plus \emph{coup}, a \textbf{beautiful blow}, and it means \emph{a lot} — so the same punch that gives Spanish its \emph{golpe} gives French its word for plenty.
\end{cousinweb}

\subsection*{Your turn}
Say these aloud:

\begin{itemize}
\raggedright
  \item they don't charge you for it all at once
  \item a knock at the door
  \item which English word is the same blow
\end{itemize}

\begin{tcolorbox}[breakable,colback=teal!4,colframe=teal!35!black,title={Before you move on}]
The blow? (\textbf{El golpe}.) All at once? (\textbf{De golpe}.) Which English word is its twin? (\textbf{Coup} — both are \emph{colaphus}.)
\end{tcolorbox}
\section[(invitar, el importe, descontar, el …]{(invitar, el importe, descontar, el golpe) — from cold}

\label{lesson:ES-R474-repaso-pago}



\begin{quote}
Four words, no prompts. What somebody does when they are paying, what the bill says, what happens to it over twelve months, and what it would have been if it had all come at once.
\end{quote}

\subsection*{You'll want to know: the four, cold}
\noindent
\begin{tabularx}{\linewidth}{@{}>{\raggedright\arraybackslash}X>{\raggedright\arraybackslash}X@{}}
\toprule
\textbf{English} & \textbf{Spanish} \\
\midrule
to invite & \textbf{invitar} \\
the amount & \textbf{el importe} \\
to deduct & \textbf{descontar} \\
the blow & \textbf{el golpe} \\
\bottomrule
\end{tabularx}

\begin{grammarlens}[title={Grammar Lens: three of the four were already in the room}]
Only one of these arrived from nowhere.

\noindent
\begin{tabularx}{\linewidth}{@{}>{\raggedright\arraybackslash}X>{\raggedright\arraybackslash}X@{}}
\toprule
\textbf{word} & \textbf{what it came out of} \\
\midrule
invitar & \textbf{el invitado}, its own participle \\
el importe & \textbf{importar}, its third sense \\
descontar & \textbf{contar}, with \emph{des-} reversing \\
\bottomrule
\end{tabularx}

\emph{El golpe} is the outsider, and even it has a relative you know in English.

So the chapter was mostly a matter of being handed things the book had been holding: \emph{invitar} had been named inside \emph{el invitado}, and \emph{importar} had told you two of its three jobs already.
\end{grammarlens}

\begin{cousinweb}
Two roots do most of the work here.

\noindent
\begin{tabularx}{\linewidth}{@{}>{\raggedright\arraybackslash}X>{\raggedright\arraybackslash}X@{}}
\toprule
\textbf{root} & \textbf{what it gave} \\
\midrule
\emph{portare}, to carry & \textbf{importe}, deporte, comportamiento \\
\emph{computare}, to reckon & \textbf{descontar}, contar, English \emph{count} \\
\bottomrule
\end{tabularx}

And \emph{golpe} is Latin \textbf{colaphus} from Greek \textbf{kólaphos}, a box on the ear — the same blow English keeps in \textbf{coup}, and the one French built \emph{beaucoup} out of. \emph{Golpe de estado} and \emph{coup d'état} are the same two words said twice.
\end{cousinweb}

\subsection*{Your turn}
Say these aloud:

\begin{itemize}
\raggedright
  \item they deduct the amount over twelve months, not all at once
  \item this one's on me
  \item which of the four did not come from a word you had
\end{itemize}

\begin{tcolorbox}[breakable,colback=teal!4,colframe=teal!35!black,title={Before you move on}]
To invite? (\textbf{Invitar}.) The amount? (\textbf{El importe}.) To deduct? (\textbf{Descontar}.) All at once? (\textbf{De golpe}.)
\end{tcolorbox}
\section[Practice]{(el curso que paga la empresa) — match it}

\label{lesson:ES-C474-sintesis-pago}



\begin{quote}
Say \textbf{el importe} and \textbf{de golpe}. The exchange below turns on whether the second applies to the first.
\end{quote}

\begin{tcolorbox}[breakable,skin=enhanced,colback=orange!6,colframe=orange!45!black,boxrule=0.5pt,arc=1mm,left=6pt,right=6pt,top=4pt,bottom=4pt,fonttitle=\bfseries,title={Read this}]
\begin{quote}
— El setenta por ciento. El resto lo pones tú, pero te lo \textbf{descuentan} en doce meses, no de \textbf{golpe}. — Ah, eso no lo sabía.
\end{quote}

And the statement you have to match to it:

\begin{quote}
Pensaba que había que pagar el \textbf{importe} completo de una vez.
\end{quote}
\end{tcolorbox}

\begin{grammarlens}[title={Grammar Lens: the match is between two phrases, not two words}]
This is a matching task, so nothing is quoted back at you. The speaker and the statement describe the same fact in different words, and you have to see it.

\noindent
\begin{tabularx}{\linewidth}{@{}>{\raggedright\arraybackslash}X>{\raggedright\arraybackslash}X@{}}
\toprule
\textbf{the speaker says} & \textbf{the statement says} \\
\midrule
no \textbf{de golpe} & \textbf{de una vez} \\
te lo descuentan en doce meses & el importe completo \\
\bottomrule
\end{tabularx}

Neither phrase appears in the other. \emph{De golpe} and \emph{de una vez} are the pair that carries the match, and if you only hold one of them the item is closed to you even though you understood every sentence.

That is worth knowing in general: a matching task rewards \textbf{paraphrase}, so the words to learn are the ones that come in pairs.
\end{grammarlens}

\subsection*{Your turn}
Say what the company pays and what happens to the rest, using \emph{descontar}.

Now say the misunderstanding twice — once with \emph{de golpe} and once with \emph{de una vez} — and hear that they are the same claim.

\begin{tcolorbox}[breakable,colback=teal!4,colframe=teal!35!black,title={Before you move on}]
What happens to the rest? (\textbf{Te lo descuentan en doce meses}.) Not how? (\textbf{De golpe}.) What did the listener think? (\textbf{Que había que pagar el importe completo de una vez}.)
\end{tcolorbox}
